\section{Πρόβλημα 3}

Θεωρούμε το max-cut game και θα βρούμε συνάρτηση δυναμικού (με απόδειξη) για τις
ακόλουθες περιπτώσεις.

\subsection{Ερώτηση (i)}

Κάθε ακμή $\{u, v\}$ έχει βάρος $w_{uv}$ που δηλώνει πόσο πολύ αντιπαθεί ο $u$ τον $v$
και αντίστροφα (συμμετρική αντιπάθεια) και κάθε παίκτης $u$ μεγιστοποιεί το
$\sum_{v \in V_u} w_{uv}$, όπου $V_u$ η μεριά του cut που δεν περιέχει τον $u$.
Θα βρούμε συνάρτηση δυναμικού με απόδειξη.

Έστω $G = (V, E)$ ο γράφος του παιγνίου με βάρη $w_{uv} \geq 0$ και έστω
$s = (s_u)_{u \in V}$ με $s_u \in \{0, 1\}$ ένα προφίλ στρατηγικών που ορίζει
διαμέριση $(S_0, S_1)$ με $S_b = \{v \in V : s_v = b\}$. Συμβολίζουμε με
$U_u = \sum_{\{u,v\} \in E,\, s_v \neq s_u} w_{uv}$ την ωφέλεια του παίκτη $u$.

Η συνάρτηση
$$
  \Phi(S) = \sum_{\substack{\{u,v\} \in E \\ s_u \neq s_v}} w_{uv}
$$
είναι ακριβής συνάρτηση δυναμικού. Θα δείξουμε ότι για κάθε παίκτη $u$ που
αλλάζει μεριά ισχύει $\Delta\Phi = \Delta U_u$.

Έστω ο παίκτης $u$ αλλάζει μεριά ενώ οι υπόλοιποι παραμένουν σταθεροί. Ορίζουμε
$N_{\mathrm{in}}(u) = \{v : \{u,v\} \in E,\, s_v = s_u\}$ τους γείτονες του $u$
στην ίδια μεριά και $N_{\mathrm{out}}(u) = \{v : \{u,v\} \in E,\, s_v \neq s_u\}$
τους γείτονες στην άλλη μεριά. Μετά την αλλαγή, οι ακμές προς $N_{\mathrm{in}}(u)$
γίνονται ακμές του cut ενώ οι ακμές προς $N_{\mathrm{out}}(u)$ παύουν να ανήκουν
σε αυτό, οπότε
$$
  \Delta\Phi = \sum_{v \in N_{\mathrm{in}}(u)} w_{uv} - \sum_{v \in N_{\mathrm{out}}(u)} w_{uv} \,.
$$
Αντίστοιχα, $U_u = \sum_{v \in N_{\mathrm{out}}(u)} w_{uv}$ πριν την αλλαγή και
$U_u' = \sum_{v \in N_{\mathrm{in}}(u)} w_{uv}$ μετά, οπότε $\Delta U_u = \Delta\Phi$.
Άρα η $\Phi$ είναι ακριβής συνάρτηση δυναμικού.

\subsection{Ερώτηση (ii)}

Κάθε παίκτης-κόμβος $v$ έχει ένα βάρος $w_v$ που δηλώνει πόσο πολύ τον αντιπαθούν
οι υπόλοιποι παίκτες (κοινή αντιπάθεια, χάνεται η συμμετρία) και κάθε παίκτης $u$
μεγιστοποιεί το $\sum_{v \in V_u} w_v$, όπου $V_u$ η μεριά του cut που δεν περιέχει
τον $u$. Θα βρούμε συνάρτηση δυναμικού με απόδειξη.

Έστω $(S_0, S_1)$ διαμέριση των κόμβων. Συμβολίζουμε με
$W_0 = \sum_{v \in S_0} w_v$ και $W_1 = \sum_{v \in S_1} w_v$ τα συνολικά βάρη
κάθε μεριάς. Η ωφέλεια του παίκτη $u \in S_b$ είναι $U_u = W_{1-b}$.

Η συνάρτηση
$$
  \Phi(S) = W_0 \cdot W_1
$$
είναι συνάρτηση δυναμικού τάξης. Θα δείξουμε ότι για κάθε παίκτη $u$ που αλλάζει
μεριά ισχύει $\mathrm{sgn}(\Delta\Phi) = \mathrm{sgn}(\Delta U_u)$.

Έστω $u \in S_0$ αλλάζει στη μεριά $S_1$. Τότε $W_0' = W_0 - w_u$ και
$W_1' = W_1 + w_u$, οπότε
$$
  \Delta\Phi = (W_0 - w_u)(W_1 + w_u) - W_0 W_1 = w_u(W_0 - W_1 - w_u) \,.
$$
Για την ωφέλεια, $U_u = W_1$ πριν και $U_u' = W_0 - w_u$ μετά, οπότε
$\Delta U_u = W_0 - w_u - W_1$. Άρα $\Delta\Phi = w_u \cdot \Delta U_u$ και αφού
$w_u > 0$, τα $\Delta\Phi$ και $\Delta U_u$ έχουν πάντα το ίδιο πρόσημο. Άρα
η $\Phi$ είναι συνάρτηση δυναμικού τάξης.
